\section{Sparse integral equations still need more than polylogarithmic separators}
\label{sec:sparse-coordinate-separators}

The explicit family of Section~\ref{sec:coordinate-separator-obstruction}
uses dense, growing integer entries.  One might hope that a sparse equation
presentation with bounded coefficients supplies the missing balanced-separator
theorem.  The following family rules out that inference too.  It has a
totally unimodular incidence presentation and a large feasible coordinate
factor after every small singleton restriction.  This remains an abstract
cone statement: neither a native triangulation realization nor a canonical
Euler function is provided.

\begin{proposition}[A sparse obstruction surviving implicit-zero deletion]
For every integer $L\ge3$, put $t=L^2$.  There is a
$(t-1)\times3t$ matrix $C$ such that:
\begin{enumerate}
\item Its entries lie in $\{-1,0,1\}$, each column has at most two nonzero
entries, each row has at most six, and $C$ is totally unimodular of rank $t-1$.
\item Its columns partition into $t$ triples summing to zero, and
$\ker C\cap\mathbb R_{\ge0}^{3t}$ has full feasible support and one
coordinate factor.
\item After restricting any $s<L/\sqrt3$ triples to arbitrary
singletons, exact implicit-zero deletion still leaves a coordinate factor
with at least
\[
                         3(L^2-s^2)>2L^2
\]
columns.
\end{enumerate}
Thus a singleton separator giving two-thirds balance requires
$s\ge\lceil\sqrt{t/3}\rceil$.  Constant coefficient size, bounded row and
column support, total unimodularity, rank $t-1$ and the local triple relations
do not imply the polylogarithmic separator hypothesis.
\end{proposition}

\begin{proof}
Take vertices $(i,j)\in(\mathbb Z/L\mathbb Z)^2$.  For each such pair use the
three directed edges
\[
 \begin{split}
 H_{ij}&:(i,j)\longrightarrow(i+1,j),\\
 D_{ij}&:(i+1,j)\longrightarrow(i,j+1),\\
 V_{ij}&:(i,j+1)\longrightarrow(i,j).
 \end{split}
\]
They form one directed triangular cycle.  For $L\ge3$ these $t$ triangles
partition all $3t$ edges of a simple undirected triangular grid on the torus.
Every vertex has degree six, with three incoming and three outgoing edges.
Let $C$ be the directed vertex-edge incidence matrix with any one row removed,
using $-1$ at a tail and $+1$ at a head.

The horizontal and vertical edges connect all vertices, so the incidence
matrix has rank $t-1$.  For completeness, every column sums to zero, and
the columns of any spanning tree become an invertible square matrix after
one row is removed: successive leaf elimination gives determinant $\pm1$.
Every square submatrix of $C$ also has determinant in $\{0,1,-1\}$.
If it has a column with at most one nonzero entry, expand along that column
and use induction.  Otherwise each column has its two opposite entries,
so its rows sum to zero and its determinant vanishes.  This proves total
unimodularity.  The three incidence columns of each directed triangle sum
to zero; summing their nonnegative cycle vectors gives the all-ones kernel
vector.  The support and encoding assertions follow immediately.

A singleton restriction deletes two edges of its triangle.  Let $I$ and $J$
be the sets of first and second indices, respectively, of the $s$ restricted
triangles.  Thus $|I|,|J|\le s$.  For every $j\notin J$ all edges of the
horizontal directed cycle in row $j$ survive; for every $i\notin I$ all edges
of the vertical directed cycle in column $i$ survive.  Call these cycles
clean, and let $G_0$ be their union.  There are at least $L-s\ge2$ clean
cycles in each direction: $s<L/\sqrt3$ and $L\ge3$ imply $L-s>1$.

The directed graph $G_0$ is strongly connected.  Each horizontal cycle meets
every vertical cycle, and a directed cycle permits travel between any two
of its vertices.  Its undirected edges also form a single graphic-matroid
component.  To see this directly, arrange the clean rows and columns in
their cyclic orders.  The rectangular cells between consecutive clean rows
and consecutive clean columns have simple cycle boundaries.  These cycles
cover all edges of $G_0$, and neighboring cells share a nonempty boundary
segment.  The cell adjacency graph is connected.  Hence their circuits join
every edge of $G_0$ into one component.

Every triangle indexed outside $I\times J$ is untouched.  If $j\notin J$,
its $H$ edge belongs to $G_0$; if $i\notin I$, its $V$ edge belongs to $G_0$.
Thus the triangle's three-column circuit meets the common component of
$G_0$ in an edge and joins all three columns to it.  Its directed cycle
indicator is also a nonnegative kernel vector positive on each of those
three coordinates.  There are at least $L^2-|I||J|\ge L^2-s^2$ such
triangles, and their edge sets are disjoint.  Let $E_*$ be the union of their
edges.  Exact implicit-zero deletion cannot remove any edge of $E_*$.

After that deletion, the remaining cone has full feasible support by the
definition of its retained coordinates.  The completeness proposition in
Section~\ref{sec:euler-coordinate-factors} identifies its finest coordinate
products with the incidence column-matroid components.  The circuits just
described survive the deletion, so one such component contains all of $E_*$.
The edge-disjoint triangle count gives
\[
                    |E_*|\ge3(L^2-s^2)>2L^2.
\]
At $s=0$ this argument includes every edge and proves the single-factor
assertion for the unrestricted cone.  Since the root has $3t=3L^2$ columns,
the displayed bound prevents two-thirds balance after every specified
restriction, completing the proof.
\end{proof}

The matrix can be generated directly from the modular edge formulas in
polynomial time.  It has at most $6t$ nonzero entries and constant-size
coefficients, with $O(t\log t)$ bits in a sparse indexed encoding.  The
circulation witnesses used in the proof are literal directed-cycle indicator
vectors.  No numerical experiment, determinant search or assumed genericity
is needed.

The restricted cone need not have full support on all its explicit columns;
the proof deliberately establishes a large surviving feasible factor instead.
It therefore covers the stronger separator mechanism that first deletes
every coordinate certified identically zero.  The lower bound is on this
particular singleton-and-coordinate-product mechanism, not on the running
time of deciding a circulation LP or another representation of these cones.

This narrows the prospective structural route without claiming a native
lower bound.  Sparsity and integral linear programming alone do not provide
the desired polylogarithmic separators.  A useful theorem must exploit
additional restrictions of authenticated knot-exterior normal equations,
their canonical Euler objectives, or a decomposition beyond literal
coordinate products.  In particular the periodic grid in this construction
is not an asserted normal-coordinate model of a knot exterior, and it says
nothing about whether a native branch could close immediately by a
nonpositive Euler certificate.  The current algorithm, its tested native
sources and its proven complexity bound are unchanged by this section.
